\section*{Bab \ref{ch:b2:corrosion} --- Korosi dan Proteksi}

\begin{solution}{exo:b2:corrosion:1}
$j = \qty{0.020}{A/m^2}$: $de/dt = 0.020 \times 65.38/(2 \times 96\,485 \times 7.13
\times 10^6) = \qty{9.5e-13}{m/s}$, yaitu \qty{0.030}{mm} per tahun.
\end{solution}

\begin{solution}{exo:b2:corrosion:2}
Baut baja pada tembaga: besi, yang kurang mulia, menjadi anode, dan
baut-baut kecil pada katode tembaga yang luas terkorosi dengan cepat. Paku
seng pada baja: seng terkorosi dan melindungi baja di sekelilingnya.
\end{solution}

\begin{solution}{exo:b2:corrosion:3}
Di celah sempit antara pelat-pelat itu dan di bawah kepala baut, tempat air
diperbarui dengan lambat dan oksigennya cepat habis: melalui \omterm{def:b2:corrosion:differential}{aerasi
diferensial}, zona-zona tertutup ini menjadi anode, dan permukaan yang
terbuka menjadi katode.
\end{solution}

\begin{solution}{exo:b2:corrosion:4}
Kaleng timah: timah, yang lebih mulia daripada besi, menjadikan baja yang
terbuka sebagai anode dari katode yang luas. Pada ember galvanis, seng
terkorosi alih-alih bajanya.
\end{solution}

\begin{solution}{exo:b2:corrosion:5}
Pengadukan menipiskan \omterm{def:b2:current-potential-curves:limiting-current}{lapisan difusi}: plato oksigen naik, kurva anodik
besi memotongnya lebih tinggi, sehingga \omterm{def:b2:corrosion:uniform}{potensial korosi} dan \omterm{def:b2:corrosion:uniform}{arus korosi}
sama-sama bertambah. Tanpa oksigen, satu-satunya reaksi katodik yang
tersisa dalam air netral adalah reduksi air yang lambat: \omterm{def:b2:corrosion:uniform}{potensial korosi}
turun dan \omterm{def:b2:corrosion:uniform}{arus korosi} menjadi sangat kecil.
\end{solution}

\begin{solution}{exo:b2:corrosion:6}
$t = 0.90 \times 5000 \times 2 \times 96\,485/(65.38 \times 0.10) = \qty{1.33e8}{s}$,
sekitar 4.2~tahun.
\end{solution}

\begin{solution}{exo:b2:corrosion:7}
Muatan per kilogram, $nF/M$: seng \qty{820}{A.h/kg}, magnesium
\qty{2.21e3}{A.h/kg}, aluminium \qty{2.98e3}{A.h/kg}. Potensial standar:
\ce{Zn^2+}/\ce{Zn} $-0.76$, \ce{Mg^2+}/\ce{Mg} $-2.36$, \ce{Al^3+}/\ce{Al}
\qty{-1.68}{V}. Di tanah yang kering dan resistif, arus harus didorong
melalui hambatan yang besar: magnesium, yang letaknya jauh di bawah besi,
menawarkan tegangan pendorong yang paling besar.
\end{solution}

\begin{solution}{exo:b2:corrosion:8}
$U = 2.0 \times 4.0 + 1.5 = \qty{9.5}{V}$; daya \qty{19}{W}; per tahun $19 \times
8766 = \qty{1.7e5}{W.h}$, \qty{1.7e2}{kW.h}.
\end{solution}

\begin{solution}{exo:b2:corrosion:9}
Di garis air, logam dibasahi air yang kaya oksigen, sebuah katode; tepat di
bawahnya, air kurang beraerasi dan zona itu menjadi anode sel yang
terbentuk bersama garis air: baja terkorosi di sana, dalam sebuah pita.
\end{solution}

\begin{solution}{exo:b2:corrosion:10}
Dalam air beraerasi, \omterm{def:b2:corrosion:uniform}{potensial korosi} baja tahan karat jatuh pada plato
pasifnya: arusnya sangat kecil, dan lapisannya memulihkan diri. Ion klorida
merusak lapisan itu pada titik-titik lemah; logam telanjang di sana
menjadi anode kecil yang berhadapan dengan katode pasif yang sangat luas,
sehingga rapat arusnya sangat besar. Di dalam sumuran, ion-ion logam
terhidrolisis dan mengasamkan larutan, klorida tertarik masuk untuk
menyeimbangkan muatan, dan \omterm{def:b2:corrosion:passivation}{lapisan pasif} tidak dapat terbentuk kembali:
sumuran itu makin dalam.
\end{solution}

\begin{solution}{exo:b2:corrosion:11}
Oksigen direduksi pada seluruh \qty{22}{cm^2}: \omterm{def:b2:current-potential-curves:ie-curve}{arus katodik} $22 \times 5 =
\qty{110}{\micro A}$, yang seluruhnya dipasok oleh oksidasi \qty{2}{cm^2}
baja di dekat sambungan, $j = \qty{55}{\micro A/cm^2}$. Kehilangannya $5.5$
kali kehilangan besi pada \qty{10}{\micro A/cm^2}, sekitar \qty{0.64}{mm}
per tahun.
\end{solution}

\begin{solution}{exo:b2:corrosion:12}
Besi imun di bawah garis \ce{Fe^2+}/\ce{Fe}, $-0.41 + 0.0296\log 10^{-6} =
\qty{-0.59}{V}$: jika dijaga di bawah sekitar \qty{-0.6}{V}, baja tidak lagi
larut. Jauh lebih rendah lagi, air direduksi pada baja dengan laju yang
makin besar (garis \ce{H+}/\ce{H2} berada pada \qty{-0.41}{V} pada pH 7 dan
\omterm{def:b2:current-potential-curves:overpotential}{overpotensial} pada baja beberapa persepuluh volt): arus terbuang dan
hidrogen mengelupas lapisan pelindung.
\end{solution}

\begin{solution}{pb:b2:corrosion:1}
\textbf{1.} \ce{Fe -> Fe^2+ + 2e-}; \ce{O2 + 2H2O + 4e- -> 4OH-}.
\textbf{2.} Reduksi oksigen, yang lambat dan mencapai plato difusinya,
membatasi \omterm{def:b2:current-potential-curves:ie-curve}{arus katodik}; kurva anodik besi memotongnya pada plato itu.
\textbf{3.} $j = \qty{0.10}{A/m^2}$: $0.10 \times 3.156 \times 10^7/(2 \times 96\,485)
\times 55.845 = \qty{913}{g}$ per meter persegi per tahun.
\textbf{4.} $913/7.87 \times 10^6 = \qty{1.16e-4}{m}$, \qty{0.12}{mm} per tahun.
\textbf{5.} $4/0.116 = 34$~tahun.
\textbf{6.} Korosinya tidak seragam: celah, perbedaan aerasi di antara
jenis-jenis tanah, dan cacat lapisan pelindung memusatkan \omterm{def:b2:current-potential-curves:ie-curve}{arus anodik} pada
daerah-daerah kecil, tempat sumuran menembus dinding jauh lebih awal.
\textbf{7.} $E^\circ = \Delta_f G^\circ/(nF)$ untuk setiap pasangan: \ce{Fe}
$-0.41$, \ce{Zn} $-0.76$, \ce{Mg} $-2.36$, \ce{Al} \qty{-1.68}{V}.
\textbf{8.} Seng, magnesium, dan aluminium, yang semuanya berada di bawah
besi.
\textbf{9.} Aluminium terpasivasi: lapisan oksidanya menghentikan
pelarutannya dan aluminium lalu hanya memberikan sedikit arus, kecuali jika
dipadukan untuk mencegah lapisan itu (yang berhasil di air laut, tidak di
dalam tanah).
\textbf{10.} Arus harus melintasi hambatan tanah: tegangan pendorong antara
anode dan baja yang dilindungi sekitar \qty{1.9}{V} untuk magnesium
dibandingkan dengan \qty{0.35}{V} untuk seng (dari potensial standar).
\textbf{11.} $3.156 \times 10^7 \times 24.305/(2 \times 96\,485 \times 0.50) =
\qty{7.95}{kg}$ per ampere per tahun.
\textbf{12.} $\pi \times 0.30 \times 10\,000 = \qty{9.42e3}{m^2}$.
\textbf{13.} \qty{94.2}{m^2}.
\textbf{14.} $0.020 \times 94.2 = \qty{1.88}{A}$.
\textbf{15.} $1.885 \times 20 \times 3.156 \times 10^7 = \qty{1.19e9}{C}$.
\textbf{16.} $1.19 \times 10^9 \times 24.305/(2 \times 96\,485 \times 0.50) =
\qty{3.00e5}{g}$, \qty{300}{kg} (atau $1.885 \times 20 \times 7.95$).
\textbf{17.} $300/15 = 20$ anode.
\textbf{18.} Hambatan tanah membatasi seberapa jauh satu anode dapat
mendorong arusnya di sepanjang pipa: anode-anode yang tersebar
melindunginya secara merata.
\textbf{19.} Sebuah penyearah mengalirkan arus dari anode inert yang
ditanam dalam hamparan anode, melalui tanah, ke dalam pipa yang dihubungkan
dengan kutub negatifnya.
\textbf{20.} $1.885 \times 5.0 + 1.5 = \qty{10.9}{V}$.
\textbf{21.} $10.9 \times 1.885 = \qty{20.6}{W}$; $20.6 \times 8766 = \qty{1.8e5}{W.h}$,
\qty{180}{kW.h} per tahun.
\textbf{22.} Air direduksi pada pipa: hidrogen terbentuk, mengelupas
lapisan pelindung dan dapat membuat baja rapuh, dan arus terbuang.
\textbf{23.} $E < -0.41 + 0.0296\log(10^{-6}) = \qty{-0.59}{V}$.
\textbf{24.} Proteksi selama dua puluh tahun memerlukan
$\boldsymbol{\approx \qty{3.0e2}{kg}}$ anode magnesium.
\end{solution}
